\section{Introduction}
\label{sec:intro}

A perception module is usually judged twice. Synthetic stress tests give free labels and
controllable factors. Small sets of real photographs give expensive labels and thin coverage. Each
judgement can mislead on its own. A synthetic test that shares a rendering lineage with the
training data reports in-distribution accuracy while looking like a robustness test. A real-photo
set that has been scored during development becomes a development set. A pipeline chosen by its
score on a small private leaderboard carries whatever that leaderboard rewarded, whether or not it
generalises. We met all three while building a chess-position recogniser. This paper is the
controlled follow-up.

Chess is a convenient testbed. Ground truth is exact: a placement FEN (Forsyth--Edwards Notation) names every one of 64
squares. Positions can be generated legally and checked for duplicates against named public pools.
Renders are cheap, so one factor can be varied while everything else is held fixed. Public sets
exist, both real photographs \cite{chessred,slcc,cvchess} and renders \cite{chesscog,cr360}, and a
strong public recogniser with released weights exists \cite{chessqueries}, so external baselines
need no retraining. The domain is narrow. What we offer are evaluation practices that were cheap
to apply here and that other sim-to-real pipelines can try, not a demonstrated transfer. The
setting is the common one: a perception module trained on renders and deployed on a camera, as in
robot manipulation, visual inspection and augmented reality.

\paragraph{Contributions.}
\begin{enumerate}
  \item A controlled render generator and a benchmark (Section~\ref{sec:benchmark}) of 20{,}000
    cell-images, 18{,}000 of them unique renders, with exact projected corners and five
    single-factor axes including the board colour convention. The design is paired: the same
    position appears at every level of a factor. Positions are disjoint from 67 inventoried
    sources, including the named public datasets and a 6.1-million-puzzle snapshot. A
    100-group partition is unscored by us, and its image hashes are committed before the
    images are released.
  \item Three documented evaluation failures (Section~\ref{sec:lies}), each stated at the strength
    its evidence supports. A leaderboard-selected pipeline scores below its own member on 16 of 18
    splits and on three further cohorts (a descriptive comparison). Two named holdouts overlapped the
    training mixes (a provenance failure). One renderer expression, the board colour convention,
    accounts for most of the per-square collapse of a state-of-the-art recogniser that had been
    read as a domain shift (a paired intervention).
  \item A matched-position, fixed-update adaptation study (Section~\ref{sec:adaptation}) with a
    preregistered real-photo retention gate. It reports what a small budget of renders changes:
    most of the zero-shot gap, and most of the convention gap when both conventions are rendered. It
    also reports what it does not resolve: the lowest elevation band reaches only 14.5 exact after
    the largest budget, and 11--23 real-photo boards are lost at every budget.
  \item A prepared release with receipts: generator, manifests, scoring and diagnostic tools, figure
    pipeline, licence register, and the test-split hash list (Section~\ref{sec:release}).
\end{enumerate}

\paragraph{What we do not claim.} No new architecture and no state of the art. Results on the
development partition are results on our own render lineage; on their own they say nothing about
real-world robustness. The public real-photo sets we report were scored during development, and we
say so where each number appears. We have not isolated the cause of the leaderboard-selected
pipeline's failure; that needs an intervention we did not run.

\paragraph{Writing assistance.} Parts of the manuscript were drafted with Codex (OpenAI) and Claude
(Anthropic) from the project's records and reports; the author verified every number against its
receipt and is responsible for the text.
